% 原书印刷xi–xii / PDF12–13
\section*{第二版前言}\phantomsection\label{front:preface-second}
\addcontentsline{toc}{section}{第二版前言}
本书的大部分引语出自阿瑟·柯南·道尔爵士，不过，若要描述这本书的历程，最贴切的也许是感恩至死乐队的那句感叹：“这一路真是漫长而奇异.”

第二版的计划大约始于六年前. 很长一段时间里，我们都在纠结应该增添什么、删去什么. 所幸，随着时间推移，统计学的发展方向日渐清晰，这些问题的答案也逐渐明朗. 我们看到，学科的趋势正在从特殊情形的优美证明，转向更复杂、更实际问题的算法解法. 这并不削弱数学与严谨性的重要性；事实上，我们发现它们变得更加重要了. 改变的是它们的应用方式.

对于熟悉第一版的读者，主要变化可简要概括如下. 渐近方法的讨论大幅扩充，独立成章. 我们更重视计算与模拟（见\hyperref[sec:5.5]{第5.5节}\footnote{译注：原文指向第5.5节；本版目录中“随机样本的生成”实际列于第5.6节.原引用保留.}及\hyperref[app:computer-algebra]{“计算机代数”附录}）；更具应用性的技术得到扩充或新增，例如自助法（bootstrapping）、EM算法、$p$值、逻辑回归与稳健回归；“杂录”和习题也新增了许多内容. 我们降低了等变性与决策论等较专门理论主题的比重，并调整了\hyperref[sec:3]{第3章}至\hyperref[sec:11]{第11章}的一些材料，使叙述更清楚.

有两点需要特别说明. 首先，关于计算机代数程序，虽然我们认为它们正成为越来越有价值的工具，但并不想强迫不持这种看法的教师使用它们. 因此，这部分以“不打扰”的方式处理：只放在附录中，并在书中一些可能有帮助的地方给出提示. 其次，我们采用了更便于查找的编号体系. 现在，定理、引理、例、定义等合用一个编号序列；例如，\hyperref[def:7.2.4]{定义7.2.4}之后是\hyperref[ex:7.2.5]{例7.2.5}，\hyperref[thm:10.1.3]{定理10.1.3}之后是\hyperref[ex:10.1.4]{例10.1.4}.

前四章只作了小幅改动. 我们重排了部分材料，特别是将不等式与恒等式分开，增添了一些例子和习题，并作了一般性的更新. \hyperref[sec:5]{第5章}也调整了顺序，把收敛一节移到较后的位置，并新增了随机变量生成一节. 第一版第7至9章中有关不变性的内容大幅缩减，合并到了\hyperref[sec:6]{第6章}；该章其余部分只作了少量编辑，主要是新增习题. \hyperref[sec:7]{第7章}扩充并更新，新增了EM算法一节. \hyperref[sec:8]{第8章}也作了少量编辑与更新，并新增$p$值一节. 在\hyperref[sec:9]{第9章}中，我们更加重视枢轴方法，因为我们意识到，“保证得到一个区间”实际上不过是“对累积分布函数作枢轴变换”. 此外，第一版第10章的决策论内容有所缩减，点估计、假设检验和区间估计在损失函数意义下的最优性分别以小节形式加入相应章节.

\hyperref[sec:10]{第10章}是全新的，力图阐明大样本推断的基础，包括Delta方法、相合性与渐近正态性、自助法、稳健估计量、得分检验等. \hyperref[sec:11]{第11章}讨论经典的单因素方差分析与线性回归；第一版中这些内容分属两章. 遗憾的是，受篇幅所限，我们删去了随机区组设计. \hyperref[sec:12]{第12章}讨论变量含误差的回归，并新增了稳健回归与逻辑回归的材料.

使用第一版授课多年之后，我们大致知道一年课程可以讲授多少内容. 例如，第二版中的下列内容应当可以在一年内讲完：
\begin{center}\begin{tabular}{ll@{\qquad}ll}
\hyperref[sec:1]{第1章} & 第1至7节 & \hyperref[sec:6]{第6章} & 第1至3节\\
\hyperref[sec:2]{第2章} & 第1至3节 & \hyperref[sec:7]{第7章} & 第1至3节\\
\hyperref[sec:3]{第3章} & 第1至6节 & \hyperref[sec:8]{第8章} & 第1至3节\\
\hyperref[sec:4]{第4章} & 第1至7节 & \hyperref[sec:9]{第9章} & 第1至3节\\
\hyperref[sec:5]{第5章} & 第1至6节 & \hyperref[sec:10]{第10章} & 第1、3、4节
\end{tabular}\end{center}
若学生开始上课时已有一定概率论基础，就可以多讲一些后面章节的材料.

最后，要感谢所有以某种方式促成第二版问世，并帮助我们改正第一版错误的人，几乎是不可能的. 感谢所有花时间寄来便条或电子邮件的学生、朋友和同事. 有些人的关键建议使我们的叙述方式发生了重大改变. 有时这些建议只是简短的便条或评论，有时则是较长的评阅意见. 有些事已过去很久，他们可能已经忘了，但我们没有忘记. 因此，感谢Arthur Cohen、David Cox爵士、Steve Samuels、Rob Strawderman和Tom Wehrly. 我们也十分感谢Jay Beder，他多年来寄来无数评论与建议，对第一版的熟悉程度可能比我们还高；还要感谢Michael Perlman及其班上的学生，就在我们写下这些话时，他们仍不断寄来意见与勘误.

这本书经历了几位编辑. 我们感谢Alex Kugashev，他在20世纪90年代中期最先建议编写第二版；感谢一直鼓励我们的编辑Carolyn Crockett. 除了我们之外，对这本书贡献最大的人也许就是第一位编辑John Kimmel，他鼓励、出版并推广了第一版. 谢谢你，John.
\begin{flushright}George Casella\\Roger L. Berger\end{flushright}

% 用户授权疑误脚注：2026-11-04
